% Transcription — fonds Alexandre Grothendieck, shelfmark 153, batch 1 (pages 1-10).
% Established from the facsimile archives/batches/153/batch-01.pdf (a mirror of
% https://grothendieck.umontpellier.fr/153.pdf), per the skill
% transcribe-grothendieck. The folder is ten pages and this is its only batch.
%
% Six of the ten leaves are transcribed. The five odd-numbered ones — 1, 3, 5,
% 7, 9 — carry his hand and his own pagination, a circled 1 to 5 at the top of
% the sheet, so the run is continuous and complete. Page 2 is the reverse of
% page 1: a computer listing, in the free right-hand third of which he wrote a
% block on the same subject as the rest; it is transcribed like any other page.
% Pages 4, 6, 8 and 10 are listings of the same run with nothing of his on
% them, and are skipped — the gaps in the page numbers are the record.
%
% The subject is the notion the folder's own title names, and the run builds it
% twice over before defining it. Page 1 is a first attempt at the axioms,
% written around a k_0-algebra Omega_0 with a distinguished element T and a
% composition (F,G) |--> F o G that is associative, has T for unit, is an
% algebra homomorphism in the first variable, and whose behaviour on G + G' and
% on GG' is set down as two equations left with nothing on their right-hand
% side; the lower half of the leaf is a jotting field on W(A), End W(A) and the
% action of Omega_0 on A. Page 2 writes the binomial polynomials
% F_n(T) = T(T-1)...(T-n+1)/n!, their two structure constants — one for the
% product F_m F_n, one for the composite F_m(F_n) — and states that the free
% abelian group they span inside Q[T] is the Z-algebra they generate. Page 3
% takes an inclusion of commutative rings k_0 in K_0 and defines Omega as the
% set of polynomials of K_0[T] taking k_0 into k_0; it is a sub-k_0-algebra
% stable under composition, it operates on k_0, and the two co-operations
% F(x+y) and F(xy) are read off K_0[X,Y] = K_0[X] tensor K_0[Y]. Page 5 is the
% definition proper — Omega a non-unitary pseudo-ring with a monoid law F o G
% subject to three axioms — with the constants Omega_0, the splitting
% Omega = Omega_0 + Omega^0, and the unitary case. Pages 7 and 9 are the
% examples: Appl(A,A) for a pseudo-ring A, End(W_k(A)), the relative Omega of a
% pair k_0 in K_0 with Z in Q as the case of interest — the binomial analyseur
% — and, named and not developed, the analyseur of divided powers and that of
% lambda-rings.
%
% The hand is drafting. The formulas, the labelled axioms and the displayed
% maps carry; the connective prose often does not and is marked rather than
% filled, page 7 being the worst of it. Two readings are flagged and not
% settled: the underlined word opening the variant of page 3, read « Lemme »;
% and the sign carried in superscript on the second of the two Delta operators
% at the foot of that page. Three things about the notation, fixed once here:
% the struck k_0<T> of page 1 is the notation Omega_0 replaces; « cts » on page
% 7 is his abbreviation for « constantes »; and the examples of page 9 are
% numbered by him in the order 4 then 3, the first numeral written over
% another.
%
% Pass: Opus 5 (claude-opus-5), 2026-09-12 — first pass, unchecked against the pages by a human.
\documentclass[11pt,a4paper]{article}
% Le préambule commun à toutes les transcriptions du fonds Grothendieck.
%
% Il tient en une idée : **l'incertitude se compose, elle ne se résout pas.**
% Une transcription qui lisse un mot douteux a détruit la seule chose pour
% laquelle on la faisait — un lecteur doit pouvoir savoir, à chaque endroit,
% ce qui était sur le feuillet et ce qui a été ajouté. Les six macros qui
% suivent sont donc l'appareil critique, et non de la décoration.
%
% Les mêmes commandes sont reconnues par `scripts/render.mjs`, qui produit la
% vue de lecture du site. Une macro ajoutée ici doit l'être aussi là-bas,
% faute de quoi le rendu échouera bruyamment — ce qui est voulu : un rendu
% silencieusement faux est pire qu'un rendu absent.

\ProvidesPackage{grothendieck}[2026/08/08 Transcriptions du fonds Grothendieck]

\usepackage{amsmath,amssymb,amsthm}
\usepackage{mathrsfs}
\usepackage{stmaryrd}
% Les diagrammes commutatifs sont partout dans ce fonds : c'est la langue
% même de la géométrie algébrique de Grothendieck, et une transcription qui
% les rendrait en prose ne transcrirait rien.
\usepackage{tikz-cd}
% Français par défaut — c'est la langue des feuillets — mais l'anglais est
% chargé aussi : la traduction et le résumé sont des documents anglais, et la
% typographie française (espaces avant les deux-points) y serait une faute.
% Ces documents basculent par \selectlanguage{english} en tête de corps.
\usepackage[english,french]{babel}
\usepackage{fontspec}
\usepackage{geometry}
\geometry{a4paper,margin=25mm,marginparwidth=28mm}
\usepackage{marginnote}
\usepackage{xcolor}
% ulem, not soul. `soul`'s \st and \ul are built on a character-by-character
% scan that XeLaTeX's font machinery breaks: the document fails with an
% undefined control sequence rather than degrading. ulem does the same two jobs
% and is Unicode-engine safe. `normalem` stops it hijacking \emph, which the
% transcriptions use for Grothendieck's own emphasis.
\usepackage[normalem]{ulem}
\usepackage{hyperref}
\hypersetup{colorlinks=true,linkcolor=black,urlcolor=black,pdfborder={0 0 0}}

% --- Métadonnées du lot -----------------------------------------------------
% Reprises telles quelles par le script de rendu, qui les affiche en tête de
% la vue de lecture. Elles fixent ce que le fichier prétend couvrir : sans
% elles, une transcription flotte sans point d'attache dans un fonds de seize
% mille feuillets.

\newcommand{\folder}[1]{\gdef\grothFolder{#1}}
\newcommand{\batch}[1]{\gdef\grothBatch{#1}}
\newcommand{\leaves}[2]{\gdef\grothFirst{#1}\gdef\grothLast{#2}}
\newcommand{\foldertitle}[1]{\gdef\grothFoldertitle{#1}}
\gdef\grothFolder{?}\gdef\grothBatch{?}\gdef\grothFirst{?}\gdef\grothLast{?}\gdef\grothFoldertitle{}

% --- Le résumé --------------------------------------------------------------
%
% Il ouvre la lecture modernisée, sous le titre « Résumé », et il est composé
% en petit. Ce n'est pas un ornement : c'est le seul passage du document qui
% ne soit pas des mathématiques, et un lecteur doit pouvoir le reconnaître
% avant de l'avoir lu — puis le sauter s'il connaît déjà le sujet, ou s'y
% arrêter s'il ne le connaît pas. Le filet qui le suit marque la reprise du
% corps.

\newenvironment{resume}
  {\small\setlength{\parskip}{0.45em}}
  {\par\medskip\hrule\medskip}

% --- Mots-clés ---------------------------------------------------------------
%
% En anglais, en clôture du résumé de la lecture modernisée : le vocabulaire
% moderne sous lequel chercher ce que les feuillets construisent. Le manifeste
% du site (`npm run manifest`) les extrait pour étiqueter la cote entière —
% cette ligne est la source unique des tags, il n'y a pas de fichier à côté.
\newcommand{\keywords}[1]{\par\smallskip\noindent
  {\footnotesize\textsc{Keywords} --- \textit{#1}}\par}

% --- L'appareil critique ----------------------------------------------------

% Le numéro de feuillet, en marge. C'est l'ancre qui permet de régler un
% désaccord sur une formule en regardant *un* feuillet plutôt qu'une chemise
% de six cents. Le site s'en sert aussi pour tourner les pages du fac-similé
% quand on fait défiler la transcription.
\newcommand{\leaf}[1]{%
  \marginnote{\footnotesize\color{gray!70}\textsf{#1}}%
  \label{leaf:#1}\ignorespaces}

% Illisible : on ne devine pas. Un mot inventé qui se lit comme les autres est
% le pire résultat possible.
\newcommand{\ill}{\textcolor{red!60!black}{[\,\ldots\,]}}

% Lecture douteuse : le mot est donné, et signalé comme incertain.
\newcommand{\uncertain}[1]{\uline{#1}}

% Ajout de l'éditeur — une lettre manquante, un mot sous-entendu.
\newcommand{\add}[1]{\textcolor{blue!50!black}{[#1]}}

% Biffé par Grothendieck lui-même. Ce qu'il a rayé fait partie du document :
% une rature dit souvent où la pensée a changé de direction.
\newcommand{\struck}[1]{\sout{#1}}

% Note du transcripteur, sur l'état matériel du feuillet.
\newcommand{\note}[1]{\par\smallskip{\small\color{gray!60!black}\textit{[#1]}}\par\smallskip}

% Note marginale de Grothendieck, distincte du corps du texte.
\newcommand{\marginal}[1]{\par\smallskip\noindent\hspace{1em}%
  {\small\color{gray!60!black}#1}\par\smallskip}

% --- Titre ------------------------------------------------------------------

\AtBeginDocument{%
  \begin{center}
    {\large\bfseries Fonds Alexandre Grothendieck --- cote n\textsuperscript{o}~\grothFolder}\\[2pt]
    {\itshape\grothFoldertitle}\\[4pt]
    {\small feuillets \grothFirst--\grothLast\ (lot \grothBatch)}\\[2pt]
    {\footnotesize\color{gray!60!black}Transcription établie d'après le fac-similé numérisé
     par l'Université de Montpellier.\\
     Les lectures incertaines sont soulignées, les illisibles notées [\,\ldots\,],
     les ajouts entre crochets.}
  \end{center}
  \vspace{1em}}

\endinput


\folder{153}
\batch{1}
\pages{1}{10}
\dating{[à partir de 1982]}
\watermark{Édition de démonstration}
\foldertitle{Analyseurs : notes manuscrites (s.d.)}

\begin{document}

\note{les cinq feuillets de mathématiques portent sa pagination propre, un
chiffre cerclé en haut du feuillet, 1 à 5, sur les pages 1, 3, 5, 7 et 9 des
archivistes. La page 2 est le verso de la page 1 : il a écrit dans le tiers
droit resté libre un bloc du même sujet, transcrit ici à son rang}

\page{1}

$k_0$ anneau de base (cas important : $k_0 = \mathbb{Z}$)

\struck{$k_0\langle T\rangle$} $k_0$-algèbre, contient $T$ élément $T$, d'où

\note{sic : le $T$ est écrit deux fois. La notation $k_0\langle T\rangle$,
biffée ici et aux deux lignes suivantes, est celle que $\Omega_0$ remplace}

$k_0[T] \overset{\varphi}{\longrightarrow}$ \struck{$k_0\langle T\rangle$}
$\Omega_0$

Opérations \quad \struck{$k_0\langle T\rangle \times k_0\langle T\rangle$}
$\longrightarrow$ \struck{$k_0\langle T\rangle$} $\Omega_0$

\note{au-dessus du produit biffé, en interligne : $\Omega_0 \times \Omega_0$}

\[
(F,G) \longmapsto F \circ G
\]

\marginal{plutôt $W(\Omega_0) \times W(\Omega_0) \to W(\Omega_0)$}

\begin{enumerate}
\item[a)] $(F \circ G) \circ H = F \circ (G \circ H)$

\item[b)] $T \circ F = F \circ T = F$ \qquad ainsi q\add{ue} $\varphi$ est
défini en termes de l'opération $F \circ G$ sur $\Omega$.

\item[c)] $\varphi(F \circ G) = \varphi(F) \circ \varphi(G)$
\note{sous le $F \circ G$ du membre de gauche, accoladé : « composition
habituelle »}

\item[d)] $F \longmapsto F \circ G : k_0\langle T\rangle \to k_0\langle T\rangle$
est \struck{\ill{}} $k_0$-alg.
\note{au-dessus de la ligne, en interligne et d'une plume qui n'appuie pas, un
mot lu \uncertain{bilatérales}}

\item[e)] $F \circ (G + G') = $

\item[f)] $F \circ (GG') = $
\end{enumerate}

\note{e) et f) sont accolées ensemble et laissées sans second membre}

$\Omega_0$ $k_0$-algèbre (év.\ sans unité)

\note{sous le $\Omega_0$, un petit signe d'appartenance couché, sans suite}

\begin{enumerate}
\item[g)] $F \circ \lambda 1 = F(\lambda)\cdot 1$ \qquad
$\bigl[$ où $\Omega_0$ opère sur $k_0$ par $(F,\lambda) \mapsto F(\lambda)$
$\bigr]$
\end{enumerate}

\[
W(\Omega_0) \times W(\Omega_0) \longrightarrow W(\Omega_0)
\]

\note{la moitié inférieure du feuillet est un champ d'essais, écrit plus vite
et sans phrases suivies ; il est donné tel quel}

\struck{$F$} \qquad $k_0$ \qquad $A$ $k_0$-algèbre commut.\ « sans unité »

$E^1$ \qquad \struck{$S$} \qquad \struck{$\operatorname{End}(M)$} \qquad
$W(A)$ \qquad $\otimes_{k_0}$-algèbre

$\operatorname{End} W(A)$ . \struck{\ill{}} structures d'algèbre
$\otimes_{k_0}$

\[
(\lambda F)(x) = \lambda F(x)
\]
\[
(F+G)(x) = F(x) + G(x)
\]
\[
(FG)(x) = F(x)G(x)
\]

structure \struck{\ill{}} d'anneau

\[
W(\Omega_0) \longrightarrow \operatorname{End} W(A)
\]

et une structure \struck{additive} \ill{}, et \ill{} à la multiplication

$\Omega_0$ opère sur $A$ \qquad $(F \circ G)(x) = F(G(x))$

\begin{tikzcd}
W(\Omega_0) \arrow[r, no head] & W(A)
\end{tikzcd}

\[
F(x+y) = \Delta F(x,y)
\qquad \Omega_0 \otimes \Omega_0
\qquad A \otimes A
\qquad
\begin{matrix} k_0 \\ | \\ \mathbb{Z} \end{matrix}
\]

\page{2}

$n \in \mathbb{N}$
\[
F_n(T) = \frac{T(T-1)\cdots(T-n+1)}{n!} \in \mathbb{Q}[T]
\]

\[
(*) \qquad F_m(T)F_n(T) = \sum_{0 \leq p \leq m+n} c^{p}_{m,n} F_p(T)
\]
\note{sous les $c^{p}_{m,n}$, accoladé : $\in \mathbb{Z}$}

\[
F_m(F_n(T)) = \sum_{0 \leq p \leq mn} d^{p}_{m,n} F_p(T)
\]

\[
\bigoplus_n \mathbb{Z}\cdot F_n \subset \mathbb{Q}[T]
\]

est la $\mathbb{Z}$-algèbre \add{comm.\ unit.} engendrée par les $F_n$
($n \in \mathbb{N}$) soumis aux relations : $F_0 = 1$ et \uncertain{ces}
(on pose alors $F_1 = T$),

\note{« comm.\ unit. » est en interligne au-dessus de « $\mathbb{Z}$-algèbre » ;
la phrase reste en suspens sur le mot souligné}

\[
F_m(P(T) + Q(T)) = \sum \gamma^{m}_{ij} F_i(P(T)) F_j(Q(T))
\]

\[
\delta : \Omega \longrightarrow \Omega \otimes_{\mathbb{Z}} \Omega
\]

\[
F_m(P(T)Q(T)) = \sum \delta^{m}_{ij} (P(T))
\]
\note{la formule s'arrête là, la parenthèse ouverte sur $F_m$ n'étant pas
refermée}

\[
F_m(x+y) \qquad F_m(xy)
\]

\[
\sum a_i T^i \qquad a_i \in \mathbb{Q}
\]

\page{3}

Soit \struck{$k_0$ anneau intègre, $K_0$ son corps des fractions} Soit
$k_0 \subset K_0$ inclusion, \uncertain{d'anneaux commutatifs}.

\[
\Omega \subset K_0[T], \qquad
\Omega = \{\, F \in K_0[T] \mid F(\lambda) \in k_0 \ \ \forall \lambda \in k_0 \,\}
\]

\begin{tikzcd}[column sep=small, row sep=small, nodes={font=\scriptsize}]
K_0[T] \arrow[r] & \operatorname{Appl}(K_0,K_0) \\
\Omega \arrow[u, hook] \arrow[r, dashed] & \operatorname{Appl}\bigl((K_0,k_0),(K_0,k_0)\bigr) \arrow[u, hook]
\end{tikzcd}

\note{le feuillet écrit les deux relations verticales comme des signes
d'inclusion, non comme des flèches}

Alors $\Omega$ est une sous-$k_0$-algèbre de $K_0[T]$, stable par
$F \circ G$.

$\Omega$ opère sur $k_0$, de façon que
\[
\begin{cases}
(F+G)(x) = F(x) + G(x) \\
(\lambda F)(x) = \lambda(F(x)) \\
(FG)(x) = F(x)G(x) \\
(F \circ G)(x) = F(G(x))
\end{cases}
\]

\note{la deuxième ligne est d'abord écrite $(\lambda F)(x) = \lambda F(x)$,
puis reprise avec les parenthèses}

\uncertain{Lemme} : si $M$ est un $k_0$-module libre, le sous-ens.\ $\Omega_M
\subset K_0[T] \otimes_{k_0} M$ formé des $F$ tels que $F(x) \in M$
$\forall x \in k_0$, s'identifie à $\Omega \otimes_{k_0} M$.

\note{le mot souligné qui ouvre cet énoncé est formé d'un $L$ suivi de quatre
jambages ; il est lu « Lemme » et n'est pas assuré}

De plus, on a\add{ura}
\[
F(x+y) \overset{?}{=} \sum_i G_i(x) H_i(y)
\]
\[
F(xy) \overset{?}{=} \sum_\alpha J_\alpha(x) K_\alpha(y)
\]

où $\sum_i G_i \otimes H_i$ est un élém.\ bien déterminé de
$\Omega \otimes_{k_0} \Omega$, et $\sum_\alpha J_\alpha \otimes K_\alpha$ est
un él.\ bien déterminé de $\Omega \otimes_{k_0} \Omega$.

\marginal{OK si $\Omega$ est un $k_0$-module \emph{libre}}

\note{cette remarque est écrite en biais, à environ 45°, le long du bord gauche
du feuillet}

En effet, dans $K_0[X,Y] = K_0[X] \otimes_{k_0} K_0[Y]$, on a\add{ura}
\[
F(X+Y) \overset{\text{déf}}{=} \Delta_a F(X,Y)
\]
\[
F(XY) \overset{\text{déf}}{=} \Delta^{\times}_{m} F(X,Y)
\]
\note{les deux opérateurs sont soulignés sur le feuillet ; le signe porté en
exposant sur le second est peu net et se lit $\times$ ou $+$, le premier n'en
portant aucun}

polynômes en $X$, $Y$, à coeff.\ $\in K_0$, qui pour des val.\ $x,y \in k_0$
\uncertain{ont} des valeurs dans $k_0$.

\[
F \circ (P+Q) = \sum_i (G_i \circ P)(H_i \circ Q)
\]
\[
F \circ (PQ) = \sum_\alpha (J_\alpha \circ P)(K_\alpha \circ Q)
\]
\[
F \circ \lambda T = \sum_\alpha J_\alpha(\lambda) K_\alpha
\qquad \Bigl( = \sum_\alpha K_\alpha(\lambda) J_\alpha \Bigr)
\]

\section*{Analyseur}

\note{le mot est de lui, souligné, en tête de la page 5}

\page{5}

\textbf{Analyseur} $\Omega$ : c'est un pseudo-anneau non unitaire (i.e.\ un
groupe additif avec multiplication associative commutative, pas néc.\
unitaire)

\[
x+y, \qquad xy
\]

avec une loi d\struck{e composition arbitraire} \add{de monoïde}
\struck{(pas néc.\ \ill{})}

\struck{loi additive}
\[
(F,G) \longmapsto F \circ G
\]

(unité notée \struck{$E$} $I$ ou $T$) pas néc.\ commutative

\note{« de monoïde » est écrit en interligne au-dessus du passage biffé}

\begin{enumerate}
\item[1)] \textbf{Pour mémoire} :
\[
\begin{cases}
+ \ \text{ass., commutative, avec } 0 \\
xy \ \text{associative, commutative, biadditive} \\
x \circ y \ \text{associative, unitaire}
\end{cases}
\]

\item[2)] $\forall G \in \Omega$, $F \longmapsto F \circ G$ est un hom.\ de
pseudo-anneaux, i.e.
\[
\begin{cases}
(F+F') \circ G = F \circ G + F' \circ G \\
(FF') \circ G = (F \circ G)(F' \circ G)
\end{cases}
\]

\item[3)] (Fonctorialité)

\item[a)] $\forall F \in \Omega$, $\exists$ famille finie
$(F'_i, F''_i)$, $F'_i, F''_i \in \Omega$, telle que
\[
F \circ (G' + G'') = \sum_i (F'_i \circ G')(F''_i \circ G'')
\]

\item[b)] $\forall F \in \Omega$, $\exists$ famille finie
$(P'_\alpha, Q''_\alpha)$
\[
F \circ (G'G'') = \sum_i (P'_\alpha \circ G')(Q''_\alpha \circ G'')
\]
\end{enumerate}

\note{a) et b) sont les deux moitiés de la condition 3) et lui sont
subordonnées sur le feuillet}

\note{en b), l'indice porté sous le signe somme est un $i$ alors que la famille
est indexée par $\alpha$}

\textbf{NB} \quad Que dire de $F \circ 0$ ? On a
\[
(F \circ 0) \circ G = F \circ \bigl(\underbrace{0 \circ G}_{0}\bigr) = F \circ 0 ,
\]
donc $F \circ 0$ est « une constante » au sens suivant $\vee$

\textbf{Déf} \quad Un $\lambda \in \Omega$ est dit « constante » si
$\lambda \circ F = \lambda$ $\forall F \in \Omega$.
$\bigl[$ Ne dépend que de la structure $F \circ G$ $\bigr]$

\textbf{Proposition} \quad L'ens.\ $\Omega_0$ des constantes de $\Omega$ est un
sous-\ill{} de $\Omega$ \ill{}, car $\lambda \circ \mu = \lambda$ pour
$\lambda, \mu \in \Omega$. Pour $F \in \Omega$, $\lambda \in \Omega_0$, on a
$F \circ \lambda \in \Omega_0$ (car
$(F \circ \lambda) \circ G = F \circ (\lambda \circ G) = F \circ \lambda$, le
$\lambda \circ G$ étant accoladé $\lambda$).

Ainsi, $\Omega$ « opère » sur $\Omega_0$ par
\[
(F,\lambda) \longmapsto F(\lambda) \overset{\text{déf}}{=} F \circ \lambda
\]
\[
\Omega^{0} = \{\, F \in \Omega \mid F \circ 0 \overset{\text{déf}}{=} F(0) = 0 \,\}
\qquad \text{On a} \quad \Omega = \Omega_0 \oplus \Omega^{0}
\]

\textbf{Déf} \quad Analyseur \textbf{unitaire} : la multiplication $xy$ admet
une unité $1$, et celle-ci est une constante, i.e.
\[
1 \circ F = 1 \qquad \forall F \in \Omega .
\]

Dans $\Omega_0$ on a \ill{} unitaire, \ill{}

\page{7}

\textbf{Ex 1} \quad Soit \struck{$k$ \ill{}} \ill{} $A$ un
\struck{\ill{}} \ill{},
\[
\Omega_A = \operatorname{Appl}(A,A) .
\]
Alors on \ill{} que $\operatorname{Appl}(E,A)$, $A$ \ill{} (\ill{} si $A$
est unitaire), et on \ill{} que \struck{$\operatorname{Appl}$}$(A,A)$, $A$
\ill{}

$\Omega$ a une structure de \ill{} ses \struck{\ill{}} $A$ est unitaire),
il y a une structure de composition, \struck{monoïdal} \ill{} unité. Les
conditions 1°) et 2°) \ill{} vérifiées. Les \struck{\ill{}} « cts » de
$\Omega$ sont les applications constantes, qui \ill{} $\Omega_0$ \ill{} :
$A$. (\ill{} unit.)

\note{« cts » est son abréviation pour « constantes ». Cette page est écrite
très vite : les formules portent, la prose qui les relie ne porte pas, et elle
est marquée plutôt que comblée}

$\Omega_A$ est \ill{} dans les $A$ l'\ill{}.

\textbf{Variante} : si $A$ est une $k$-algèbre, on peut prendre
($\Omega_0 \simeq A$),
\[
\underline{\Omega}_{A/k} = \operatorname{End}(W_k(A))
\]
\ill{} plus général, si $A$ est une \ill{} \ill{} types, on dispose \ill{}
Anneau commutatif \uncertain{unitaire} \ill{}

\[
\Omega(\struck{\ill{}}) = \operatorname{End}_{\text{faisceaux d'ens.}}(A) ,
\]
\ill{} deux structures $xy$ et $\lambda \circ y$ \ill{} (Unitaire si $A$
l'est)

\[
\boxed{(\Omega_A)_0 \simeq \Gamma A \, ?}
\]

Dans le cas où $A = k$, on trouve $\Omega \simeq k[T]$, et \ill{} la
\add{addition,} multiplication \ill{} la bilat.\ des \ill{}

\ill{} analyseurs \struck{de tous les « standards »} des polynômes. Dans ce
cas, les conditions 3°) \ill{} vérifiées. Anneau des constantes est $k$.

On peut \ill{} prendre $k[[T]]^{+}$, \ill{} produit et $\lambda \circ y$,
\ill{} réduit à zéro (ou les $a_0 + a_1T + \cdots$ avec $a_0$
\uncertain{nilpotent})

Soit $\Omega$ un analyseur, \struck{arbitraire} \ill{}, $\Omega_0$ l'ens.\
des cts, alors
\[
F \longmapsto (\lambda \mapsto F \circ \lambda) : \Omega \longrightarrow
\operatorname{Appl}(\Omega_0,\Omega_0) = \Omega_{\Omega_0}
\]
est un « hom.\ d'analyseurs », induisant \ill{} sur les \struck{\ill{}}
constantes. On est content \ill{} que $\Omega \to$
\struck{$\operatorname{App}$} $\Omega_{\Omega_0}$ soit injectif !

\page{9}

\textbf{Ex 2} \quad Relatif : un couple $k_0 \subset K_0$ d'anneaux
commutatifs,
\[
\Omega \subset K_0[T], \qquad
\Omega = \{\, P \in K_0[T] \mid P(\lambda) \in k_0 \ \ \forall \lambda \in k_0 \,\}
\]
structures $+$, $\cdot$, $\circ$ induites par celles de $K_0[T]$. Le cas le
plus intéressant $\to$ \ill{} celui des corps
($\mathbb{Z} \subset \mathbb{Q}$ :

\textbf{L'analyseur binomial} \qquad (unitaire)

\textbf{Ex 4} \quad Analyseur des puissances divisées \qquad (pas unitaire)

\textbf{Ex 3} \quad Analyseur des $\lambda$-anneaux \qquad (unitaire).

\note{un trait sépare l'Ex 2 des deux derniers ; ceux-ci sont écrits dans
l'ordre 4 puis 3, le premier chiffre tracé par-dessus un autre, et une flèche
partant de leur hauteur remonte vers ce trait}

\end{document}
